\Needspace{60mm}
\section{分年度表现与稳健性}
\subsection{分年度表现}
\begingroup\small\setlength{\tabcolsep}{4pt}
\begin{longtable}{lrrrrrrr}
\caption*{\textbf{表C3　全部 A 股分年度表现}}\\
\toprule
年份 & 月数 & \makecell{多空\\收益} & \makecell{多空月\\胜率} & \makecell{Q5\\收益} & \makecell{内部基准\\收益} & \makecell{多头\\超额} & 平均IC \\ \midrule
\endfirsthead
\caption*{表C3　全部 A 股分年度表现（续）}\\\toprule
年份 & 月数 & \makecell{多空\\收益} & \makecell{多空月\\胜率} & \makecell{Q5\\收益} & \makecell{内部基准\\收益} & \makecell{多头\\超额} & 平均IC \\ \midrule
\endhead
\midrule\multicolumn{8}{r}{续下页}\\\endfoot
\bottomrule\endlastfoot
2007* & 11 & 37.39 & 63.64 & 184.68 & 150.96 & 33.71 & 0.0896 \\
2008 & 12 & 12.36 & 66.67 & -53.13 & -56.76 & 3.63 & 0.0650 \\
2009 & 12 & 30.89 & 83.33 & 164.05 & 135.40 & 28.66 & 0.0918 \\
2010 & 12 & -8.44 & 50.00 & 6.87 & 14.86 & -7.99 & 0.0038 \\
2011 & 12 & 12.38 & 58.33 & -26.52 & -33.32 & 6.80 & 0.0523 \\
2012 & 12 & 6.39 & 41.67 & 10.94 & 6.09 & 4.84 & 0.0431 \\
2013 & 12 & -4.65 & 25.00 & 17.91 & 26.76 & -8.85 & 0.0090 \\
2014 & 12 & 38.12 & 75.00 & 74.61 & 47.88 & 26.73 & 0.0993 \\
2015 & 12 & 12.71 & 50.00 & 91.24 & 82.89 & 8.35 & 0.0212 \\
2016 & 12 & 28.65 & 75.00 & 17.53 & 3.19 & 14.33 & 0.0952 \\
2017 & 12 & 22.51 & 66.67 & 0.72 & -11.80 & 12.51 & 0.0723 \\
2018 & 12 & 10.22 & 58.33 & -28.71 & -30.16 & 1.45 & 0.0526 \\
2019 & 12 & 2.37 & 66.67 & 28.81 & 30.54 & -1.73 & 0.0369 \\
2020 & 12 & -2.70 & 50.00 & 16.54 & 18.97 & -2.44 & 0.0192 \\
2021 & 12 & 11.54 & 50.00 & 32.26 & 23.75 & 8.51 & 0.0645 \\
2022 & 12 & 11.76 & 66.67 & -1.22 & -10.78 & 9.56 & 0.0602 \\
2023 & 12 & 25.24 & 66.67 & 18.04 & 8.82 & 9.22 & 0.0813 \\
2024 & 12 & 13.13 & 66.67 & 10.51 & 1.94 & 8.56 & 0.0659 \\
2025 & 12 & -0.17 & 50.00 & 37.61 & 43.00 & -5.39 & 0.0495 \\
2026* & 1 & 1.80 & 100.00 & 3.83 & 3.52 & 0.31 & 0.1173 \\
\end{longtable}
\noindent\footnotesize 收益与胜率单位为\%，多头超额单位为百分点。年度按持有起始月归属；*为部分年度，报告实际累计收益，不向全年外推。胜率分母为该行月数。\par
\endgroup

2007年仅11期、2026年仅1期，为实际累计收益。18个完整年度中，多空收益为正14/18年，复利口径的多头超额为正13/18年，两条年度序列的相关系数为0.94，总体同步；多空不利年份为2010、2013、2020、2025。按第3.3节的判据，空头端驱动的年份为2015、2018和2019年：三年的多头部分为$-$0.21、2.09、$-$1.41个百分点，空头部分为14.42、8.12、3.95个百分点；2023年空头部分占64\%，接近该界限。2015年复利多头超额为8.35个百分点，算术多头部分为$-$0.21个百分点，差异来自逐月求和与复利运算的不同，而非仅由基准上涨放大造成。2025年多空接近零（$-$0.17\%），两端算术贡献4.32与$-$4.18个百分点相抵，但不能直接相加为年度复利收益；同年平均Rank IC为0.0495，排序关联仍在，最高层却跑输基准。

本报告的稳健性证据来自三方面：三个股票池方向一致（第4.1、4.2节）；三个日历块的平均Rank IC均为正（第4.3节）；五层与十层的年化净收益排序一致（第4.5节）。去极值方法敏感性见第5.2节，IC衰减与样本外检验留待后续研究。

% BEGIN WINSOR SENSITIVITY 20261029
\Needspace{85mm}
\subsection{去极值方法敏感性}
\begingroup\small\setlength{\tabcolsep}{3pt}
\begin{table}[H]\centering
\caption*{\textbf{表S1　去极值方法敏感性}}
\begin{tabular*}{\textwidth}{@{\extracolsep{\fill}}lrrrrr@{}}\toprule
去极值方案 & \makecell{Rank IC\\均值} & ICIR & \makecell{回归因子收益\\均值（\%/月）} & \makecell{多空年化\\净收益（\%）} & \makecell{Q5算术年化超额\\（百分点）} \\ \midrule
MAD5 & 0.0566 & 0.550 & 0.40 & 13.00 & 5.31 \\
MAD3 & 0.0565 & 0.522 & 0.40 & 12.34 & 4.80 \\
1\%/99\%分位数缩尾 & 0.0563 & 0.563 & 0.38 & 13.32 & 5.55 \\
\bottomrule\end{tabular*}
\par\smallskip\raggedright\footnotesize 全部 A 股，228期；仅改变去极值方法，其余设定及输入S样本不变。ICIR未年化；多空为扣除双腿成本后的CAGR；Q5超额为相对原内部等权毛收益基准的算术年化超额。分位点采用线性插值。差值按未舍入结果计算。
\end{table}
\endgroup

本节完成去极值敏感性分析，比较MAD5、MAD3与1\%/99\%分位数缩尾；IC衰减与样本外检验仍留待后续。表S1中，三种方案的平均回归因子收益均为正（0.38\%至0.40\%/月），平均Rank IC为0.0563至0.0566，方向一致；ICIR为0.522至0.563。相对MAD5，MAD3与分位数方案的多空年化净收益分别变化$-$0.66、+0.32个百分点，Q5算术年化超额分别变化$-$0.51、+0.25个百分点。三种方案五层年化净收益均随层号递增。按可加总的算术年化口径，MAD3毛多空下降0.55个百分点，双腿成本增加0.01个百分点，净收益下降主要来自毛收益变化。由此，原结论的正向关联与分层单调性得到支持，但收益幅度仍依赖去极值选择，不能称为完全稳健。

% END WINSOR SENSITIVITY 20261029
